\documentclass[11pt, a4paper]{article}

% --- PAQUETES PRINCIPALES ---
\usepackage[utf8]{inputenc}
\usepackage[spanish, es-tabla, es-noquoting]{babel}
\usepackage[margin=2.5cm]{geometry}
\usepackage{amsmath, amssymb, mathtools}
\usepackage{siunitx}
\usepackage{graphicx}
\usepackage{booktabs}
\usepackage{tabularx}
\usepackage[most]{tcolorbox}
\usepackage{listings}
\usepackage{xcolor}

% --- PAQUETES PARA GRÁFICAS Y CIRCUITOS ---
\usepackage{pgfplots}
\usepackage[american]{circuitikz}
\usetikzlibrary{shapes.geometric, arrows, positioning, babel}
\pgfplotsset{compat=1.18}

% --- DEFINICIÓN DE COLORES ---
\definecolor{codegreen}{rgb}{0,0.6,0}
\definecolor{codegray}{rgb}{0.5,0.5,0.5}
\definecolor{codepurple}{rgb}{0.58,0,0.82}
\definecolor{backcolour}{rgb}{0.96,0.96,0.96}

% --- CONFIGURACIÓN DE CÓDIGO (PYTHON) ---
\lstset{
    backgroundcolor=\color{backcolour},   
    commentstyle=\color{codegreen},
    keywordstyle=\color{magenta},
    numberstyle=\tiny\color{codegray},
    stringstyle=\color{codepurple},
    basicstyle=\ttfamily\footnotesize,
    breaklines=true,                 
    numbers=left,                    
    numbersep=5pt,                  
    language=Python,
    literate={á}{{\'a}}1 {é}{{\'e}}1 {í}{{\'i}}1 {ó}{{\'o}}1 {ú}{{\'u}}1
             {Á}{{\'A}}1 {É}{{\'E}}1 {Í}{{\'I}}1 {Ó}{{\'O}}1 {Ú}{{\'U}}1
             {ñ}{{\~n}}1 {Ñ}{{\~N}}1
}

% --- CONFIGURACIÓN DE SIUNITX ---
\sisetup{output-decimal-marker = {,}, per-mode = symbol, inter-unit-product = \ensuremath{{}\cdot{}}}

% --- CABECERAS ---
\usepackage{fancyhdr}
\pagestyle{fancy}
\fancyhf{}
\lhead{\textbf{Circuitos Electrónicos III (IMT-221)}}
\rhead{Sesión 4.2: Transitorios Inductivos y Flyback}
\cfoot{\thepage}
\renewcommand{\headrulewidth}{0.4pt}

% --- ENTORNOS PERSONALIZADOS ---
\newtcolorbox{aviso}[1][]{colback=red!5!white, colframe=red!75!black, fonttitle=\bfseries, title=#1}
\newtcolorbox{definicion}[1][]{colback=blue!5!white, colframe=blue!75!black, fonttitle=\bfseries, title=#1}

\begin{document}

\begin{center}
    {\LARGE \textbf{Apuntes de Cátedra: Sesión 4.2}}\\[0.3cm]
    {\large \textbf{Supresión de Transitorios Inductivos: Física, Modelado y Protección ADC}}\\
    \textit{Docente: M.Sc. Ing. Bernardo Quiroga Turdera}
    \vspace{0.3cm}
    \hrule
\end{center}

\vspace{0.3cm}

\section*{Objetivos de Aprendizaje}
\begin{itemize}
    \item Comprender la física matemática de los transitorios inductivos (conmutación en cargas reactivas).
    \item Modelar y resolver analíticamente las respuestas temporales del circuito con y sin diodo Flyback.
    \item Conectar la supresión de picos inductivos con la seguridad del módulo analógico del Hito 1.
\end{itemize}

% ==========================================
% SECCIÓN 1: FÍSICA DE LA INTERRUPCIÓN
% ==========================================
\section{La Física de la Interrupción Inductiva (El Conflicto Matemático)}
Explicación de la contradicción física al apagar un motor o bobina. Ecuación Gobernadora:
\begin{equation}
    v_L(t) = L \frac{di(t)}{dt}
\end{equation}

\textbf{Análisis Conceptual:}
\begin{itemize}
    \item Un inductor almacena energía en su campo magnético y se opone a los cambios bruscos de corriente (principio de inercia electromagnética).
    \item Cuando un transistor de control (conmutador) pasa del estado de conducción (ON) al estado de corte (OFF), la corriente que circulaba de forma estable por la bobina ($I_{init}$) es obligada a extinguirse en un tiempo extremadamente corto ($\Delta t \to 0$).
    \item Dado que $\frac{di}{dt}$ tiende a un valor negativo gigantesco, el inductor genera un voltaje inducido inverso masivo (\textit{inductive kickback}), cuya polaridad es opuesta a la de su alimentación ordinaria.
    \item \textbf{Consecuencia física:} La tensión en la terminal de conmutación del transistor se eleva a miles de voltios, provocando la avalancha por sobretensión y la inmediata destrucción destructiva del semiconductor de potencia.
\end{itemize}

% ==========================================
% SECCIÓN 2: ESCENARIO CRÍTICO SIN DIODO
% ==========================================
\section{Análisis del Escenario Crítico SIN Diodo Flyback}
\textbf{El Circuito Equivalente Real (RLC Underdamped):} Ningún interruptor es perfecto; al abrirse, el colector o drenador del transistor posee una pequeña capacitancia parásita hacia tierra ($C_p \approx \qty{100}{\pico\farad}$).

\begin{figure}[h!]
    \centering
    \begin{circuitikz}[american, scale=0.9, transform shape]
        \draw (0,0) node[vcc]{$V_{CC}$ (\qty{12}{\volt})} to[L, l=$L$ (Motor)] (0,-2) to[R, l=$R$ (Devanado)] (0,-4);
        \draw (0,-4) to[short, -*] (2,-4) node[right]{Nodo de Conmutación ($v_{node}$)};
        \draw (0,-4) to[cute open switch, l=Transistor] (0,-7) node[ground]{};
        \draw (2,-4) to[C, l=$C_p$ ($\qty{100}{\pico\farad}$)] (2,-7) node[ground]{};
    \end{circuitikz}
    \caption{Esquema Eléctrico Equivalente en el Apagado SIN Diodo Flyback.}
\end{figure}

\textbf{Ecuación Diferencial del Sistema:}\\
Al abrirse el interruptor, la corriente del inductor se desvía a la malla formada por la resistencia del devanado, la bobina y la capacitancia parásita:
\begin{equation}
    L \frac{d^2q(t)}{dt^2} + R_L \frac{dq(t)}{dt} + \frac{q(t)}{C_p} = V_{CC}
\end{equation}

\textbf{Solución para Tensión de Nodo ($v_{node}(t) = q(t)/C_p$):}\\
Dado que $C_p$ es extremadamente pequeño, el factor de amortiguamiento $\alpha = \frac{R}{2L}$ es despreciable frente a la frecuencia de resonancia natural $\omega_0 = \frac{1}{\sqrt{L C_p}}$. El sistema es altamente subamortiguado:
\begin{equation}
    v_{node}(t) \approx V_{CC} - V_{CC} e^{-\alpha t} \cos(\omega_d t) + I_{init} \sqrt{\frac{L}{C_p}} e^{-\alpha t} \sin(\omega_d t)
\end{equation}

\textbf{El Pico de Tensión Teórico (Worst-Case):}\\
El término del seno domina por completo la amplitud máxima:
\begin{equation}
    V_{peak} \approx V_{CC} + I_{init} \sqrt{\frac{L}{C_p}}
\end{equation}
\textbf{Ejemplo Práctico:} Para $V_{CC} = \qty{12}{\volt}$, $L = \qty{10}{\milli\henry}$, $R_L = \qty{5}{\ohm}$ y $C_p = \qty{100}{\pico\farad}$, la corriente inicial es $I_{init} = \qty{2.4}{\ampere}$. El pico teórico es:
\begin{equation}
    V_{peak} \approx 12\text{ V} + 2.4\text{ A} \times \sqrt{\frac{10^{-2}\text{ H}}{10^{-10}\text{ F}}} = 12\text{ V} + 2.4 \times 10,000 = \mathbf{24,012\text{ V}}
\end{equation}
\textit{(Este es el fenómeno que se muestra en el panel izquierdo de la gráfica interactiva flyback\_transient\_analysis.png en Studio, donde el voltaje se eleva a la escala de kilovoltios en menos de un microsegundo).}

\newpage
% ==========================================
% SECCIÓN 3: SOLUCIÓN DE INGENIERÍA
% ==========================================
\section{Solución de Ingeniería: El Diodo Flyback (Rueda Libre)}

\begin{figure}[h!]
    \centering
    \begin{circuitikz}[american, scale=0.9, transform shape]
        \draw (0,0) node[vcc]{$V_{CC}$ (\qty{12}{\volt})} -- (2,0);
        \draw (0,0) to[L, l=$L$] (0,-2) to[R, l=$R$] (0,-4);
        \draw (2,0) to[D, l=$D_1$ (Flyback 1N4007), invert] (2,-4);
        \draw (0,-4) to[short, -*] (3,-4) node[right]{Nodo de Conmutación ($v_{node}$)};
        \draw (0,-4) -- (2,-4);
        \draw (1.5,-4) to[cute open switch, l=Transistor (Abriéndose $t=0$)] (1.5,-7) node[ground]{};
    \end{circuitikz}
    \caption{Esquema Eléctrico con Protección Flyback.}
\end{figure}

\textbf{Ecuaciones de Control del Clamping:}
\begin{itemize}
    \item \textbf{Intervalo 1 (Transistor Cerrado - ON):} El nodo de conmutación está a $\approx \qty{0}{\volt}$. El cátodo del diodo está a $\qty{12}{\volt}$ y el ánodo a $\qty{0}{\volt}$. El diodo está bloqueado (OFF).
    \item \textbf{Intervalo 2 (Transistor Abriéndose - Transitorio):} El transistor se corta. La tensión del colector se eleva rápidamente por encima de $V_{CC}$. En cuanto alcanza la tensión de sujeción:
    \begin{equation}
        v_{node}(t) = V_{CC} + V_{on} \approx \qty{12.7}{\volt}
    \end{equation}
    El diodo se polariza en directa (ON), recortando instantáneamente cualquier pico superior y salvaguardando al transistor.
\end{itemize}

\subsection*{Modelo Matemático del Desgaste Exponencial (ON de Rueda Libre)}
\textbf{La Malla de Descarga:} Con el diodo encendido, la corriente de la bobina circula en un lazo cerrado recirculando por el devanado resistivo y el diodo:
\begin{equation}
    L \frac{di(t)}{dt} + R_L i(t) + V_{on} = 0 \implies \frac{di(t)}{dt} + \frac{R_L}{L}i(t) = -\frac{V_{on}}{L}
\end{equation}
\textbf{Solución Analítica de la Corriente:}
\begin{equation}
    i(t) = \left( I_{init} + \frac{V_{on}}{R_L} \right) \exp\left( -\frac{t}{\tau} \right) - \frac{V_{on}}{R_L}
\end{equation}
Donde la constante de tiempo térmica de descarga es $\tau = \frac{L}{R_L}$.

\textbf{Tiempo de Extinción de Corriente ($t_{off}$):}\\
Dado que el diodo solo permite conducción unidireccional, la corriente se extingue exponencialmente hasta llegar exactamente a cero. Evaluando $i(t_{off}) = 0$:
\begin{equation}
    t_{off} = \tau \ln\left( \frac{V_{CC} + V_{on}}{V_{on}} \right)
\end{equation}
\textbf{Cálculo de Diseño:} Para $\tau = \frac{\qty{10}{\milli\henry}}{\qty{5}{\ohm}} = \qty{2}{\milli\second}$, y $V_{CC} = \qty{12}{\volt}$:
\begin{equation}
    t_{off} = 2\text{ ms} \times \ln\left( \frac{12\text{ V} + 0.7\text{ V}}{0.7\text{ V}} \right) \approx 2\text{ ms} \times 2.899 = \mathbf{5.798\text{ ms}}
\end{equation}
\textit{(Esta dinámica se representa detalladamente en el panel derecho de flyback\_transient\_analysis.png, donde la corriente de descarga disminuye suavemente de 2.4A a 0A en 5.8 ms, mientras que el voltaje del nodo de conmutación se mantiene seguro a 12.7V).}

\newpage
% ==========================================
% SECCIÓN 4: CONEXIÓN PFI
% ==========================================
\section{Conexión Práctica con el Hito 1 del PFI}
Para asegurar que los estudiantes comprendan la relevancia de este circuito, debes guiar la clase hacia la arquitectura del módulo de sujeción analógica de su proyecto.

\vspace{0.4cm}
\begin{center}
\begin{tikzpicture}[node distance=3.2cm, auto, thick, scale=0.8, transform shape]
    \node[draw, rectangle, fill=gray!20, text width=3cm, align=center, minimum height=1.5cm] (motor) {\textbf{Carga Inductiva (Motor)}\\(Protegida por 1N4007)};
    \node[draw, rectangle, fill=blue!15, text width=3cm, align=center, minimum height=1.5cm, right of=motor, node distance=3.8cm] (shunt) {\textbf{Resistencia Shunt / Entrada}\\(Sujeción Negativa 1N4148)};
    \node[draw, rectangle, fill=orange!15, text width=3cm, align=center, minimum height=1.5cm, right of=shunt, node distance=4cm] (zener) {\textbf{Divisor $R_s$ y Zener}\\(Sujeción Positiva 5.1V)};
    \node[draw, rectangle, fill=green!15, text width=3cm, align=center, minimum height=1.5cm, right of=zener, node distance=3.8cm] (adc) {\textbf{Entrada ADC (Micro)}\\(Rango Seguro 0-5V)};

    \draw[->, >=stealth, line width=1.5pt] (motor) -- (shunt);
    \draw[->, >=stealth, line width=1.5pt] (shunt) -- (zener);
    \draw[->, >=stealth, line width=1.5pt] (zener) -- (adc);
\end{tikzpicture}
\end{center}
\vspace{0.4cm}

\textbf{La Planta de Conmutación (Protección Primaria):}\\
En su circuito, el motorreductor amarillo de prueba (3-12V) genera transitorios inductivos masivos al apagarse o al modular su velocidad mediante PWM de alta frecuencia. La colocación del diodo 1N4007 en paralelo directo a las bornas del motor extingue los picos de conmutación en la fuente.

\textbf{Protección Complementaria del ADC (Fijación de Piso y Techo):}\\
Incluso con el Flyback activo, el "ruido de tierra" y las inductancias parásitas de los cables largos de protoboard pueden transferir picos transitorios rápidos a la línea de adquisición analógica del microcontrolador.
\begin{itemize}
    \item El \textbf{Diodo Zener BZX55C5V1} (5.1V) sujeta cualquier sobrepico positivo por encima de $\qty{5.1}{\volt}$ (Protección de Techo).
    \item El \textbf{Diodo rápido de señal 1N4148} (con un tiempo de recuperación extremadamente veloz de $t_{rr} \le \qty{4}{\nano\second}$) actúa como un limitador ultra rápido que drena voltajes negativos espurios inferiores a $\qty{-0.3}{\volt}$ directamente a GND, protegiendo las compuertas lógicas de silicio de la entrada analógica del microcontrolador.
\end{itemize}

% ==========================================
% SECCIÓN 5: PYTHON CODE
% ==========================================
\section{Código Python para el Taller Computacional}
Este es el script completo para generar la simulación (Pilar 1 - Modelado Analítico de la rúbrica).

\begin{lstlisting}
# ==============================================================================
# UNIVERSIDAD CATÓLICA BOLIVIANA "SAN PABLO" - SEDE TARIJA
# Sesión 4.2: Modelado del Transitorio Inductivo y Clamping por Diodo Flyback
# Docente: M.Sc. Ing. Bernardo Quiroga Turdera
# ==============================================================================
import numpy as np
import matplotlib.pyplot as plt
import seaborn as sns

# 1. PARÁMETROS DEL SISTEMA (Motor y Diodo Flyback)
Vcc = 12.0        # V (Tensión de alimentación)
L = 10.0e-3       # H (Inductancia nominal del devanado de motor: 10 mH)
R = 5.0           # Ohm (Resistencia interna de devanado: 5 Ohm)
I_init = Vcc / R  # A (Corriente en estado estacionario antes del corte: 2.4 A)
V_on = 0.7        # V (Caída de tensión directa del diodo 1N4007)
Cp = 100e-12      # F (Capacitancia parásita enbornes del switch: 100 pF)

# 2. ESCENARIO SIN DIODO (Oscilación RLC Subamortiguada - Escala Microsegundos)
omega_0 = 1.0 / np.sqrt(L * Cp)
alpha = R / (2.0 * L)
omega_d = np.sqrt(omega_0**2 - alpha**2)

t_us = np.linspace(0, 15, 1000)   # Vector de tiempo en us
t_s_us = t_us * 1e-6              # Conversión a segundos

# Ecuaciones analíticas para carga parásita del switch sin protección
A = -Cp * Vcc
B = (I_init - alpha * A) / omega_d
q = Cp * Vcc + np.exp(-alpha * t_s_us) * (A * np.cos(omega_d * t_s_us) + B * np.sin(omega_d * t_s_us))
v_node_no_diode = q / Cp
i_inductor_no_diode = np.exp(-alpha * t_s_us) * (
    (I_init * np.cos(omega_d * t_s_us)) +
    (((Vcc - alpha * I_init * L) / (L * omega_d)) * np.sin(omega_d * t_s_us))
)

# 3. ESCENARIO CON DIODO FLYBACK (Fijación de Voltaje y Descarga RL - Milisegundos)
t_ms = np.linspace(0, 10, 1000)   # Vector de tiempo en ms
t_s_ms = t_ms * 1e-3              # Conversión a segundos

tau = L / R
t_off = tau * np.log((Vcc + V_on) / V_on)  # Momento exacto donde la corriente se extingue

i_inductor_with_diode = np.zeros_like(t_s_ms)
v_node_with_diode = np.zeros_like(t_s_ms)

for idx, t in enumerate(t_s_ms):
    if t < t_off:
        i_inductor_with_diode[idx] = (I_init + V_on/R) * np.exp(-t/tau) - V_on/R
        v_node_with_diode[idx] = Vcc + V_on
    else:
        i_inductor_with_diode[idx] = 0.0
        v_node_with_diode[idx] = Vcc

# 4. CONFIGURACIÓN DEL ENTORNO GRÁFICO (Seaborn para Calidad de Publicación)
sns.set_theme(style='whitegrid', palette='colorblind', font='DejaVu Sans')
fig, axes = plt.subplots(2, 2, figsize=(16, 10))

# Título de Insight
fig.suptitle('Análisis Comparativo del Transitorio Inductivo: Con vs. Sin Diodo Flyback\n'
             'Demostración de sujeción de sobretensión destructiva para Rango de Operación Seguro',
             fontsize=16, fontweight='bold', y=0.98)

# Panel Voltaje SIN diodo
axes[0, 0].plot(t_us, v_node_no_diode / 1e3, color='crimson', linewidth=2)
axes[0, 0].set_title('Tensión en el Switch SIN Diodo Flyback\n(Pico Destructivo en us)', fontsize=12, fontweight='bold')
axes[0, 0].set_ylabel('Voltaje del Nodo ($kV$)')
axes[0, 0].set_xlabel('Tiempo ($\mu s$)')

# Panel Corriente SIN diodo
axes[0, 1].plot(t_us, i_inductor_no_diode, color='orange', linewidth=2)
axes[0, 1].set_title('Corriente de Bobina SIN Diodo Flyback (Interrupción Rápida)', fontsize=12, fontweight='bold')
axes[0, 1].set_ylabel('Corriente ($A$)')
axes[0, 1].set_xlabel('Tiempo ($\mu s$)')

# Panel Voltaje CON diodo
axes[1, 0].plot(t_ms, v_node_with_diode, color='teal', linewidth=2)
axes[1, 0].set_title('Tensión en el Switch CON Diodo Flyback\n(Clamping Seguro a 12.7V en ms)', fontsize=12, fontweight='bold')
axes[1, 0].set_ylabel('Voltaje del Nodo ($V$)')
axes[1, 0].set_xlabel('Tiempo ($ms$)')
axes[1, 0].axhline(y=Vcc+V_on, color='purple', linestyle='--', label='Clamping: Vcc + Von (12.7V)')
axes[1, 0].legend()

# Panel Corriente CON diodo
axes[1, 1].plot(t_ms, i_inductor_with_diode, color='royalblue', linewidth=2)
axes[1, 1].set_title('Corriente de Bobina CON Diodo Flyback (Descarga Exponencial)', fontsize=12, fontweight='bold')
axes[1, 1].set_ylabel('Corriente ($A$)')
axes[1, 1].set_xlabel('Tiempo ($ms$)')

plt.tight_layout(pad=2.0)
plt.show()
\end{lstlisting}

\begin{aviso}[Consejo para la Clase Práctica (Dinámica ABET)]
Cuando los estudiantes presenten la defensa de este Hito, pídeles que realicen en vivo una modificación en el script cambiando la inductancia $L$ de $\qty{10}{\milli\henry}$ a $\qty{50}{\milli\henry}$. La teoría y el script predecirán de inmediato que la corriente tardará más en extinguirse (más tiempo de descarga de rueda libre) y ellos deberán comprobar visualmente en su osciloscopio físico si este incremento en la constante de tiempo $\tau$ coincide con el modelo y el cálculo de RMSE.
\end{aviso}

\end{document}
